\section{Secretaría}

Secretaría es quien crea y formaliza la matrícula. Todo lo demás se apoya en
ese trabajo.

\subsection{Buscar un estudiante}

\begin{enumerate}[leftmargin=*,itemsep=0.2em]
\item Entre a \textbf{Estudiantes}.
\item Escriba el nombre o el RUN (con o sin puntos) y pulse \textbf{Enter}, o
  use los filtros de año, curso y estado.
\item Pulse \textbf{Ver ficha} en la fila del estudiante.
\end{enumerate}

\captura{04-estudiantes}{0.95}{Listado de estudiantes con filtros por año, curso y estado.}

La marca \textbf{Solo con pendientes} deja a la vista los expedientes que
todavía tienen documentos por entregar o por validar.

\subsection{La ficha del estudiante}

La ficha se abre en pestañas. \textbf{Salud}, \textbf{PIE} y
\textbf{Documentos} muestran solo lo que su rol puede ver.

\captura{05-ficha-general}{0.95}{Ficha del estudiante: folio, versión y corrección con motivo.}

\begin{itemize}[leftmargin=*,itemsep=0.2em]
\item \textbf{General}: datos personales, dirección, matrículas del estudiante
  y versiones de la ficha.
\item \textbf{Autorizaciones}: imagen, convivencia y paracetamol, con firma y
  fecha.
\item \textbf{Documentos}: el expediente, con el estado de cada requisito.
\item \textbf{Salud} y \textbf{PIE}: lo que el apoderado declaró y el programa
  de integración.
\end{itemize}

\subsection{Documentos del expediente}

\captura{06-ficha-expediente}{0.95}{Expediente documental: cada requisito se marca recibido, no aplica o pendiente.}

\begin{enumerate}[leftmargin=*,itemsep=0.2em]
\item Abra la pestaña \textbf{Documentos}.
\item Elija el requisito y suba el archivo: PDF, JPG o PNG de hasta 15 MB.
\item El portal comprueba que el archivo sea realmente del tipo declarado y
  guarda una huella (SHA-256) para detectar cambios.
\item Marque el requisito como \textbf{recibido}, \textbf{no aplica} o
  \textbf{pendiente}. Ningún requisito bloquea la matrícula.
\end{enumerate}

\subsection{Autorizaciones}

\captura{07-ficha-autorizaciones}{0.95}{Autorizaciones con su estado y la opción de revocar indicando el motivo.}

Cada autorización conserva quién la firmó y cuándo. Para dejarla sin efecto se
pulsa \textbf{Revocar} y se escribe el motivo; el portal guarda el cambio y la
autorización deja de contar como vigente. Las que pasan su fecha cambian solas
a \textbf{vencida}: el portal revisa los vencimientos cada 15 minutos.

\subsection{Crear o renovar una matrícula}

\begin{enumerate}[leftmargin=*,itemsep=0.2em]
\item Entre a \textbf{Matrícula}.
\item Elija el tipo de documento (RUN o IPE) y escriba el número.
\item Pulse \textbf{Buscar}. Si el estudiante ya existe, el portal ofrece
  \textbf{Renovar} con los datos del año anterior; si no existe, ofrece
  \textbf{Nueva}.
\end{enumerate}

\captura{08-asistente-busqueda}{0.9}{Búsqueda previa a la matrícula: RUN o IPE del estudiante.}

\captura{09-asistente-paso1}{0.9}{Primer paso del asistente: identificación del estudiante.}

El asistente tiene ocho pasos:

\begin{enumerate}[leftmargin=*,itemsep=0.2em]
\item \textbf{A · Estudiante}: identificación, curso y datos personales.
\item \textbf{B · Familia}: madre, padre, apoderado principal y secundario.
\item \textbf{C · Emergencia y retiro}: contactos y personas autorizadas.
\item \textbf{D · Salud}: situación médica, alergias y salida al baño.
\item \textbf{E · PIE}: participación en el programa de integración.
\item \textbf{F · Imagen}: autorización de uso de imagen.
\item \textbf{G · Convivencia}: compromiso de convivencia escolar.
\item \textbf{Revisión}: la ficha completa con el PDF a la vista.
\end{enumerate}

Cada paso se guarda al pulsar \textbf{Continuar}; puede volver atrás sin
perder lo escrito. En la revisión final se firma y se confirma con una
declaración expresa.

\subsection{Corregir una ficha}

Si un dato está equivocado después de firmar:

\begin{enumerate}[leftmargin=*,itemsep=0.2em]
\item Abra la ficha y pulse \textbf{Corregir ficha (nueva versión)}.
\item Escriba el motivo de la corrección.
\item Corrija los datos y cierre. Se crea una \textbf{versión nueva}: la
  anterior queda como reemplazada y no se modifica.
\end{enumerate}

\subsection{Anular una matrícula}

La anulación exige motivo y deja la matrícula sin efecto, pero conserva todo
su historial. Después de anular se puede iniciar una matrícula nueva para el
mismo estudiante y el mismo año.
